% !TEX root = ../nc-quadric-classification.tex
% 2026-11-21 / v1.35
\section{一般の幾何的 cubic 族の相対内在的実現}
\label{sec:relative-cubic-families}

本節では、可換二次曲面を中心とするという制限を外し、
普遍 Noetherian 性と canonical hull を相対 tails の認識に結び付ける。
Serre 平方関係は、まず指定された周期の独立な入力として明記する。
次節の Picard 構成が非線型の普遍族上でこの入力を供給する。

\subsection{周期を選ぶ前の canonical compact 圏}

$R$ を可換 Noetherian $k$-代数とし、
$C$ を定理~\ref{thm:relative-geometric-cubic-noetherian} の仮定を満たす
正の向きの $R$-線形 $\mathbb Z$-代数とする。
右表現可能加群を $P_i=e_iC$、$\mathcal O_i=\pi P_i$ と書く。
三次関係を $R_i$、四次 overlap を $W_i$ とする。

\begin{proposition}[一般 geometric cubic 族の canonical tails]
\label{prop:general-cubic-canonical-tails}
任意の可換 Noetherian $R$-代数 $R'$ 上で
$\operatorname{QGr}(C_{R'})$ は locally Noetherian であり、
任意の連続四項の直和はその非有界導来圏の compact tilting generator である。
特に
\[
 T=\bigoplus_{i=0}^3\mathcal O_{-i},\qquad
 \Gamma=\operatorname{End}(T)
\]
と置くと、$\Gamma$ は階数 $24$ の有限 projective $R$-代数であり、
\[
 D(\operatorname{QGr}C)\simeq D(\operatorname{Mod}\Gamma),\qquad
 D(\operatorname{QGr}C)^c\simeq\operatorname{Perf}\Gamma.
 \tag{GC-Morita}\label{eq:general-cubic-canonical-morita}
\]
この比較は実際の module-object 引戻しの導来関手と整合する。
compact 圏は相対 smooth proper で、相対 Serre 関手とその任意基底変更を持つ。
この命題では二周期の存在を仮定しない。
\end{proposition}

\begin{proof}
完全な証明は相対モジュライ論文 \cite[\texttt{relative_cubic_families}]{QuadricModuli2026} に置く。
\end{proof}

\subsection{実際の二周期からの全相対条件}

$C$ に二周期同型 $\beta:C\to C(2)$ を指定し、
\[
 B_n=C_{0,2n},\qquad a\star b=a\,\beta^m(b)
 \quad(a\in B_m,\ b\in B_n)
 \tag{GC-prod}\label{eq:general-cubic-period-product}
\]
と定める。$\check B\simeq C^{\mathrm{ev}}$ の下で、
$B$ の次数移動 $(1)$ は $C$ の right tails 上で
$\mathcal O_i\mapsto\mathcal O_{i-2}$ という完全自己同値 $s$ を与える。
矢も $\beta^{-1}$ で移す。単に対象の置換を指定しているのではない。
式~\eqref{eq:general-cubic-canonical-morita} によって
$s$ を表す可逆 perfect $\Gamma$ 双加群を $P_s$ と書く。

\begin{theorem}[Serre 適合二周期族の相対効果性]
\label{thm:geometric-cubic-relative-effectivity}
上の $C,\beta$ に対し、基底の fppf 局所で
\[
 P_s\otimes_\Gamma^{\mathbf L}P_s
       \simeq\mathcal S_\Gamma^{-1}[2]
 \tag{GC-root}\label{eq:general-cubic-serre-input}
\]
が与えられ、任意の基底変更ではこの同型を引き戻すとする。
このとき有限型 $R$-基底上の系
\[
 \mathcal C_{R'}=\operatorname{qgr}(B\otimes_RR'),\qquad
 \mathcal O_{R'}=\pi(B\otimes_RR')
 \tag{GC-family}\label{eq:general-cubic-relative-family}
\]
は定義~\ref{def:relative-abelian-system}、
\ref{def:relative-h2} の全条件を満たす。
特に $\operatorname{Spec}R$ が元の局所有限型 $k$-site に属すれば、
実際の $\mathfrak H_2(R)$ の対象を与える。
平方根、平方同型、二周期、cubic helix および AS4 被覆は、
構成に用いる witness であり、出力の標識には加えない。

各 affine Noetherian 基底変更上で全相対 $\chi_q$ も成立する。
$R$ が元の site に属さない場合は、
有限型 $R$-test 上で同じ全公理が成立するという意味に限る。
\end{theorem}

\begin{proof}
完全な証明は相対モジュライ論文 \cite[\texttt{relative_cubic_families}]{QuadricModuli2026} に置く。
\end{proof}

\begin{remark}[Hom の向きで表示する場合]
\label{rem:general-cubic-right-orientation}
framed 五つ組を $D_{ij}=\operatorname{Hom}(E_i,E_j)$ として
回復する場合には、右加群のために
\[
 C_{ij}=D_{-j,-i},\qquad \beta_C=\beta_D^{-1}
\]
を用いる。このとき $E_i$ は $\pi e_{-i}C$ に対応し、
四項 foundation の頂点・全 Hom・積が一致する。
canonical 四周期も反転し、$\beta_D^2=\theta_D$ が
頂点スカラーを除いて成立すれば、$C$ 側でも対応する平方関係が成立する。
相対 Serre の比較は $C$ の canonical compact 圏で行う。
反対の積を暗黙に同じ環と同一視しない。
\end{remark}
